\title[Israel 3: Jueces--David]{\bfseries{Israel, el pueblo de Dios}}
\subtitle{Estudio 3. De los Jueces a David}
\author{}\institute{}\date{}
\titlegraphic{\includegraphics[width=0.85\linewidth,height=0.38\textheight,keepaspectratio]{portada.jpg}\\[-1pt]
  {\fontsize{5}{6}\selectfont\color{sepia}\textit{Nota.} Elaboración propia: 1~S 16:7 en texto hebreo consonántico (BHS).}}
\begin{document}
\frame{\titlepage}

% ===================================================
\section{Introducción}
\begin{frame}{¿Cansados de ser distintos?}
\begin{columns}[T,onlytextwidth]
\begin{column}{0.50\textwidth}
\begin{cita}Vosotros me seréis un reino de sacerdotes y una nación santa.\fuente{(Éx 19:6)}\end{cita}
\begin{cita}Constitúyenos ahora un rey que nos juzgue, como tienen todas las naciones.\fuente{(1~S 8:5)}\end{cita}
\begin{cita}No ha desechado Dios a su pueblo.\fuente{(Ro 11:2; cf. 1~S 12:22)}\end{cita}
\end{column}
\begin{column}{0.46\textwidth}\small
\destacar{Una ruina interna.} La vocación se quiebra \emph{desde dentro}: levitas, sacerdotes y tribus hacen «lo recto ante sus ojos» (Dt 12:8; Jue 21:25).\par\medskip
\destacar{Una promesa en pie.} Libertadores, un remanente, una moabita en el linaje de Judá y el anuncio de un rey \emph{ungido} (1~S 2:10).\par\medskip
\begin{notalam}[Sobre las citas]
Traducción propia desde el texto hebreo (BHS) y griego (NA28). Cuando la numeración difiere, se da la española y entre corchetes la hebrea.
\end{notalam}
\end{column}
\end{columns}
\end{frame}

\begin{frame}{Los ejes de la serie en esta etapa}
\centering
\begin{tikzpicture}
 \foreach \x/\n/\t/\d/\r [count=\ovc,evaluate=\ovc as \ovk using int(\ovc+0)] in {0/1/{Nación santa}/{¿Cómo se quiebra la vocación sacerdotal y quiénes la preservan?}/{Jue 2:11--19; 17--21; Rt 1:16; 1~S 2; 8:5--7; 12:23},
   4.2/2/{Restauración}/{¿Qué esperanza nace en medio de la ruina?}/{Jue 2:16--18; Rt 4:14--17; 1~S 2:10, 35; 12:22; 16:13},
   8.4/3/{Proyección}/{¿Cómo releen Isaías, Jeremías, Oseas y Pablo esta etapa?}/{Is 1:26; 11:1; Jer 23:5--6; Os 3:4--5; Ro 11:1--5}}{\begin{scope}[visible on=<\ovk->]
  \node[caja,anchor=north] (b\n) at (\x,0) {\parbox[t][3.0cm][t]{3.55cm}{\RaggedRight
    {\small\textbf{\textcolor{ladrillo}{\t}}}\\[5pt]\small \d\par\vfill{\scriptsize\color{sepia}(\r)}}};
  \node[circle,fill=cafe,text=piedra,font=\scriptsize\bfseries,inner sep=1.6pt] at (b\n.north west) {\n};\end{scope}}
 \begin{scope}[visible on=<3->]\draw[flecha] (b1.east) -- (b2.west);\end{scope}
 \begin{scope}[visible on=<3->]\draw[flecha] (b2.east) -- (b3.west);\end{scope}
\end{tikzpicture}
\elabpropia
\vspace{6pt}
\begin{clave}
\textbf{Uso de «Israel».} En esta etapa «Israel» designa a todo el pueblo; la distinción Israel\,/\,Judá asoma al final (1~S 11:8; 18:16; 2~S 2:4, 9--10).
\end{clave}
\end{frame}

% ===================================================
\section{Cronología}
\begin{frame}{Secuencia de la etapa}
\centering
\begin{tikzpicture}[x=0.76cm]
\draw[line width=1.6pt,cafe,-{Latex[length=3mm]}] (0,0) -- (15.8,0);
\foreach \x/\t/\b [count=\ovc,evaluate=\ovc as \ovk using int(\ovc+0)] in {0.5/{Muerte de Josué}/{Jue 2:8},
                      2.35/{Débora y Barac}/{Jue 4--5},
                      4.2/{Gedeón}/{Jue 6--8},
                      6.05/{Jefté}/{Jue 11--12},
                      7.9/{Sansón}/{Jue 13--16},
                      9.75/{Elí, Samuel, Afec}/{1~S 1--7},
                      11.6/{Saúl rey}/{1~S 9--11},
                      13.45/{David ungido}/{1~S 16},
                      15.2/{Gilboa}/{1~S 31}}{\begin{scope}[visible on=<\ovk->]
  \fill[oro] (\x,0) circle (4pt);
  \node[above,font=\scriptsize\bfseries,text=ladrillo,text width=1.3cm,align=center] at (\x,0.2) {\t};
  \node[below,font=\tiny,text=sepia,text width=1.3cm,align=center] at (\x,-0.2) {\b};\end{scope}}
\begin{scope}[visible on=<9->]\draw[decorate,decoration={brace,amplitude=5pt,mirror},sepia] (0.3,-0.85) -- (8.2,-0.85) node[midway,below=5pt,font=\scriptsize\itshape,align=center]{Jueces (c. 1200--1050 a.\,e.\,c.)\\Rut «en los días de los jueces» (Rt 1:1)};\end{scope}
\begin{scope}[visible on=<9->]\draw[decorate,decoration={brace,amplitude=5pt,mirror},sepia] (9.5,-0.85) -- (15.4,-0.85) node[midway,below=5pt,font=\scriptsize\itshape,align=center]{Transición a la monarquía\\(c. 1050--1010 a.\,e.\,c.)};\end{scope}
\end{tikzpicture}
\notafig{Elaboración propia. Secuencia narrativa; las fechas absolutas son aproximadas y están en debate.}
\vspace{8pt}
\begin{clave}
\textbf{Marco.} De la muerte de Josué (Jue 2:8) a la unción de David en Hebrón (2~S 2:4): la Edad del Hierro~I, un Israel de aldeas sin capital, sin rey y sin ejército permanente.
\end{clave}
\end{frame}

\begin{frame}{¿Cuánto duró el período?}
\begin{columns}[T,onlytextwidth]
\begin{column}{0.42\textwidth}\footnotesize
\begin{itemize}
\item \textbf{1~R 6:1}: 480 años del Éxodo al 4.º año de Salomón.
\item \textbf{Suma de Jueces}: c. 390 años en sucesión (cálculo propio); sumando Elí, Saúl y David, la cifra se excede.
\item Explicación habitual: \destacar{jueces regionales superpuestos}.
\item Seder Olam encaja las cifras en los 480 años (Guggenheimer, 1998).
\end{itemize}
\end{column}
\begin{column}{0.54\textwidth}
\centering
\begin{tikzpicture}[x=1cm,y=1cm]
 \draw[dashed,sepia,line width=0.6pt] (5.4,0.55) -- (5.4,-3.15);
 \node[etiq,text=sepia,anchor=north] at (5.4,-3.15) {límite de 1~R 6:1};
 \node[anchor=west,font=\tiny\bfseries,text=ladrillo] at (0,0.42) {1~R 6:1};
 \fill[ladrillo] (0,-0.17) rectangle (5.4,0.17);
 \begin{scope}[visible on=<2->]\node[font=\tiny,text=white] at (2.7,0) {480 años: del Éxodo al 4.º año de Salomón};\end{scope}
 \begin{scope}[visible on=<2->]\node[anchor=west,font=\tiny\bfseries,text=ladrillo] at (0,-0.62) {Suma de Jueces en sucesión};\end{scope}
 \begin{scope}[visible on=<3->]\fill[oro] (0,-1.21) rectangle (4.39,-0.87);\end{scope}
 \begin{scope}[visible on=<3->]\node[font=\tiny,text=cafeosc] at (2.2,-1.04) {c. 390 años};\end{scope}
 \begin{scope}[visible on=<4->]\draw[oro!90!ladrillo,dashed,fill=oro!18!white] (4.39,-1.21) rectangle (6.5,-0.87);\end{scope}
 \begin{scope}[visible on=<4->]\node[font=\tiny,text=cafeosc] at (5.45,-1.04) {+ Elí, Saúl, David};\end{scope}
 \begin{scope}[visible on=<4->]\node[anchor=west,font=\tiny\bfseries,text=ladrillo] at (0,-1.66) {Jueces regionales y simultáneos (Jue 10:7)};\end{scope}
 \begin{scope}[visible on=<5->]\fill[celeste!55!white,draw=ladrillo,line width=0.4pt] (0.6,-2.22) rectangle (3.3,-1.92);\end{scope}
 \begin{scope}[visible on=<5->]\node[font=\tiny,text=cafeosc] at (1.95,-2.07) {Amón, al este};\end{scope}
 \begin{scope}[visible on=<6->]\fill[celeste!55!white,draw=ladrillo,line width=0.4pt] (1.9,-2.67) rectangle (4.6,-2.37);\end{scope}
 \begin{scope}[visible on=<6->]\node[font=\tiny,text=cafeosc] at (3.25,-2.52) {filisteos, al oeste};\end{scope}
 \begin{scope}[visible on=<6->]\draw[decorate,decoration={brace,amplitude=3pt,mirror},oro!90!ladrillo,line width=0.7pt] (1.9,-2.8) -- (3.3,-2.8)
   node[midway,below=3pt,etiq]{a la vez};\end{scope}
\end{tikzpicture}
\notafig{Elaboración propia a partir de 1~R 6:1 y Jue 10:7; proporciones aproximadas.}
\end{column}
\end{columns}
\vspace{6pt}
\begin{clave}
Las cifras de Jueces no se pensaron como una sola cadena: si varias opresiones y liberaciones fueron regionales, los 480 años de 1~R 6:1 pueden contenerlas.
\end{clave}
\end{frame}

\begin{frame}{Dos problemas textuales}
\footnotesize\renewcommand{\arraystretch}{1.25}
\begin{tabular}{>{\RaggedRight\arraybackslash}p{1.9cm}>{\RaggedRight\arraybackslash}p{10.1cm}}
\toprule
{\color{ladrillo}\bfseries Texto} & {\color{ladrillo}\bfseries Problema}\\\midrule
Hch 13:20 & NA28: 450 años \emph{antes} de los jueces (Egipto, desierto, conquista). Texto bizantino (RV): 450 años \emph{de} jueces.\\
1~S 13:1 & TM: Saúl «tenía un año» al comenzar a reinar y «reinó dos años»: señal casi segura de cifras perdidas en la copia. Hch 13:21 le atribuye 40 años.\\
\bottomrule
\end{tabular}
\elabpropia
\vspace{8pt}
\begin{columns}[T,onlytextwidth]
\begin{column}{0.48\textwidth}
\begin{notalam}[Fechas convencionales]
Saúl c.~1025--1010; David desde c.~1010 a.\,e.\,c. (2~S 5:4--5). Son estimaciones; la duración del reinado de Saúl no puede fijarse con el texto.
\end{notalam}
\end{column}
\begin{column}{0.48\textwidth}
\begin{notalam}[Crítica textual]
Hch 13:20 tiene dos formas en los manuscritos: el texto crítico (NA28) y el bizantino ordenan de modo distinto los 450 años. Reina-Valera sigue la forma bizantina.
\end{notalam}
\end{column}
\end{columns}
\end{frame}

% ===================================================
\section{Arqueología}
\begin{frame}{El escenario de la etapa}
\begin{columns}[c,onlytextwidth]
\begin{column}{0.46\textwidth}\centering\includegraphics[width=\linewidth,height=0.78\textheight,keepaspectratio]{mapa_jueces.pdf}\par\vspace{1pt}{\fontsize{4.5}{5.5}\selectfont\color{sepia}\textit{Nota.} Elaboración propia (GSHHS/WDB, Basemap); ubicaciones aproximadas.\par}\end{column}
\begin{column}{0.50\textwidth}\scriptsize\RaggedRight
{\scriptsize\RaggedRight
\tikz[baseline=-0.6ex]\draw[oro,line width=1.5pt,line cap=round,dash pattern=on 3.4pt off 1.8pt](0,0)--(0.6,0);~Elimelec y Noemí a Moab (Rt 1:1)\par
\tikz[baseline=-0.6ex]\draw[ladrillo,line width=1.5pt,line cap=round](0,0)--(0.6,0);~Noemí y Rut regresan a Belén (Rt 1:22)\par
\tikz[baseline=-0.6ex]\draw[cafe,line width=1.5pt,line cap=round,dash pattern=on 1.2pt off 1.4pt](0,0)--(0.6,0);~El arca: Eben-ezer $\to$ Asdod… $\to$ Quiriat-jearim (1 S 4–7)\par}
\medskip

\textbf{\textcolor{cafe}{Sierra central}}: Silo, Mizpa, Guibeá; las aldeas israelitas del Hierro~I.\par\smallskip
\textbf{\textcolor{cafe}{Pentápolis filistea}}: Gaza, Asdod, Ascalón, Gat, Ecrón (Jue 3:3).\par\smallskip
\textbf{\textcolor{cafe}{Frontera}}: Bet-semes y el valle de Elá (Qeiyafa).\par\smallskip
\textbf{\textcolor{cafe}{Norte}}: Tabor (Jue 4), Harod (Jue 7), Gilboa y Bet-seán (1~S 31).\par\smallskip
\textbf{\textcolor{cafe}{Rut}}: de Belén a Moab y de regreso, cruzando el Jordán junto a Jericó.\par\medskip
\precision{Eben-ezer y otros sitios tienen identificación debatida.}
\end{column}
\end{columns}
\end{frame}

\begin{frame}{Aldeas y filisteos}
\centering
\begin{tikzpicture}
 \foreach \x/\n/\t/\d/\r in {0/1/{Aldeas de la sierra}/{Más de 250 aldeas sin muralla en el Hierro~I; casa de cuatro espacios (\emph{bet av}).}/{Finkelstein, 1988; Stager, 1985},
   3.15/2/{Medinet Habu}/{Ramsés~III combate a los Pueblos del Mar, entre ellos los \emph{Peleset} (\heb{פְּלִשְׁתִּים}).}/{c. 1177 a.\,e.\,c.; Dothan, 1982},
   6.3/3/{Ecrón}/{«Akish, hijo de Padi, rey de Ecrón»: el nombre de Aquis de Gat (1~S 27:2). Siglo \textsc{vii}.}/{Gitin et~al., 1997},
   9.45/4/{Gat}/{Ostracón con \emph{’LWT} y \emph{WLT}, nombres cercanos a «Goliat». No es Goliat.}/{Maeir et~al., 2008}}{
  \node[caja,anchor=north] (b\n) at (\x,0) {\parbox[t][3.0cm][t]{2.65cm}{\RaggedRight
    {\small\textbf{\textcolor{ladrillo}{\t}}}\\[4pt]\scriptsize \d\par\vfill{\tiny\color{sepia}(\r)}}};
  \node[circle,fill=cafe,text=piedra,font=\scriptsize\bfseries,inner sep=1.6pt] at (b\n.north west) {\n};}
\end{tikzpicture}
\notafig{Elaboración propia a partir de las fuentes citadas en cada recuadro.}
\vspace{4pt}
\begin{notalam}[Debate arqueológico: los huesos de cerdo]
En las aldeas de la sierra casi no aparecen huesos de cerdo, frecuentes en las ciudades filisteas. Faust (2006) lo lee como marcador de identidad coherente con Lv 11:7; Hesse y Wapnish (1997) advierten que la cría de cerdos depende también del clima y la economía. La cuestión sigue abierta.
\end{notalam}
\end{frame}

\begin{frame}{Silo, Bet-semes y Qeiyafa}
\begin{columns}[T,onlytextwidth]
\begin{column}{0.42\textwidth}\footnotesize
\begin{itemize}
\item \textbf{Silo}: incendio a mediados del s. \textsc{xi} a.\,e.\,c. (Finkelstein et~al., 1993). ¿Tras Afec? Samuel calla; Jer 7:12 y Sal 78:60 lo recuerdan.
\item \textbf{Bet-semes}: frontera con Ecrón, adonde vuelve el arca (1~S 6).
\item \textbf{\emph{Pim}} (1~S 13:21): palabra oscura aclarada por pesas inscritas \emph{pym}.
\end{itemize}
\end{column}
\begin{column}{0.55\textwidth}
\centering
\begin{tikzpicture}
 \node[bloqueazul,text width=4.2cm] (q) at (0,0) {Khirbet Qeiyafa\\[-1pt]{\tiny\mdseries fortaleza sobre el valle de Elá, c.~1020--980 a.\,e.\,c.}};
 \begin{scope}[visible on=<2->]\node[bloque,text width=1.75cm,font=\tiny,anchor=north] (a) at (-2.35,-0.95) {\textbf{Judá de David}\\¿Saaraim, 1~S 17:52?\\[1pt]{\color{sepia}Garfinkel}};\end{scope}
 \begin{scope}[visible on=<3->]\node[bloque,text width=1.75cm,font=\tiny,anchor=north] (b) at (0,-0.95) {\textbf{Entidad del norte}\\quizá la de Saúl\\[1pt]{\color{sepia}Finkelstein y Fantalkin}};\end{scope}
 \begin{scope}[visible on=<4->]\node[bloqueoro,text width=1.75cm,font=\tiny,anchor=north] (c) at (2.35,-0.95) {\textbf{Cananea o filistea}\\[1pt]{\color{sepia}otros autores}};\end{scope}
 \begin{scope}[visible on=<4->]\foreach \n in {a,b,c}{\draw[flecha] (q.south) -- (\n.north);}\end{scope}
\end{tikzpicture}
\notafig{Elaboración propia a partir de Garfinkel et~al. (2016) y Finkelstein y Fantalkin (2012).}
\end{column}
\end{columns}
\vspace{6pt}
\begin{notalam}[Debate: la «cronología baja»]
La «cronología baja» de Finkelstein rebaja la datación de varios estratos y reduce el reino de David; Mazar (1990) y otros la rechazan. Qeiyafa es una de las piezas centrales de esa discusión.
\end{notalam}
\end{frame}

% ===================================================
\section{Jueces y Rut}
\begin{frame}{El ciclo: una espiral}
\begin{columns}[c,onlytextwidth]
\begin{column}{0.50\textwidth}
\centering
\begin{tikzpicture}[every node/.style={font=\scriptsize}]
\foreach \a/\t/\r [count=\i,evaluate=\i as \ovk using int(\i+0)] in {90/{Apostasía}/{2:11--13},18/{Opresión}/{2:14--15},-54/{Clamor}/{3:9; 6:6},-126/{Libertador}/{2:16, 18},162/{Reposo y muerte}/{3:11; 2:19}}{\begin{scope}[visible on=<\ovk->]
  \node[draw=cafe,line width=0.8pt,fill=piedra,rounded corners=3pt,text width=1.75cm,align=center,inner sep=2.5pt] (c\i) at ({2.4*cos(\a)},{1.75*sin(\a)}) {\textbf{\textcolor{ladrillo}{\t}}\\[-1pt]{\tiny\color{sepia}Jue \r}};\end{scope}}
\begin{scope}[visible on=<5->]\foreach \a/\b in {1/2,2/3,3/4,4/5,5/1}{\draw[-{Latex[length=2mm]},oro!90!ladrillo,line width=1.1pt] (c\a) to[bend left=18] (c\b);}\end{scope}
\begin{scope}[visible on=<5->]\node[text width=1.9cm,align=center,font=\tiny\itshape,text=sepia] at (0,-0.05) {«más que sus padres» (2:19)};\end{scope}
\end{tikzpicture}
\elabpropia
\end{column}
\begin{column}{0.46\textwidth}
\begin{cita}[fontupper=\scriptsize\itshape]YHWH levantó jueces que los libraron [\ldots]. Pero cuando moría el juez, volvían a corromperse más que sus padres.\fuente{(Jue 2:16, 19)}\end{cita}
\begin{notalam}[Filología: \heb{שֹׁפֵט} (\emph{shofet})]
No es un magistrado de tribunal, sino un líder carismático que libra y gobierna. El texto insiste en que YHWH lo «levanta»: el verdadero Juez es Él (11:27).
\end{notalam}
\begin{notalam}[Filología: el clamor]
\heb{וַיִּזְעֲקוּ} (\emph{va-yiz'aqú}, 3:9): Dios responde a un gemido, no a un arrepentimiento pleno; ahí está la semilla de la restauración.
\end{notalam}
\end{column}
\end{columns}
\end{frame}

\begin{frame}{Los jueces mayores}
\centering\footnotesize\renewcommand{\arraystretch}{1.25}
\begin{tabular}{>{\RaggedRight\arraybackslash}p{1.75cm}>{\RaggedRight\arraybackslash}p{1.65cm}>{\RaggedRight\arraybackslash}p{2.1cm}>{\RaggedRight\arraybackslash}p{1.35cm}>{\RaggedRight\arraybackslash}p{3.75cm}}
\toprule
{\color{ladrillo}\bfseries Juez} & {\color{ladrillo}\bfseries Tribu} & {\color{ladrillo}\bfseries Opresor} & {\color{ladrillo}\bfseries Texto} & {\color{ladrillo}\bfseries Rasgo}\\\midrule
Otoniel & Judá & Cusán-risataim & 3:7--11 & Modelo inicial; reposo de 40 años\\
Aod & Benjamín & Eglón de Moab & 3:12--30 & Libertador zurdo; reposo de 80 años\\
Débora y Barac & Efraín\,/ Neftalí & Jabín y Sísara & 4--5 & Profetisa y jueza; Jael completa la victoria\\
Gedeón & Manasés & Madián & 6--8 & Rechaza la corona; su efod es tropiezo\\
Jefté & Galaad & Amón & 10:6--12:7 & Voto temerario; guerra contra Efraín\\
Sansón & Dan & Filisteos & 13--16 & Nazareo; vive «según sus ojos»\\
\bottomrule
\end{tabular}
\elabpropia
\end{frame}

\begin{frame}{Una espiral descendente}
\centering
\begin{tikzpicture}[x=1cm,y=1cm]
 \foreach \i/\t/\d/\s [count=\ovc,evaluate=\ovc as \ovk using int(\ovc+0)] in {0/Otoniel/{«reposó la tierra»\\3:11}/bloque,
   1/Aod/{«reposó la tierra»\\3:30}/bloque,
   2/Débora/{«reposó la tierra»\\5:31}/bloque,
   3/Gedeón/{«reposó»; pero su efod es tropiezo\\8:27--28}/bloque,
   4/Jefté/{Israel mata a Israel\\12:6}/bloqueoro,
   5/Sansón/{«recta a mis ojos»\\14:3}/bloqueoro,
   6/{Jue 17--21}/{sin rey; guerra civil\\Jue 20}/bloqueazul}{\begin{scope}[visible on=<\ovk->]
   \node[\s,text width=1.45cm,font=\tiny,anchor=north] (j\i) at ({\i*1.82},{-\i*0.5}) {{\scriptsize\bfseries \t}\\[1pt]\d};\end{scope}}
 \begin{scope}[visible on=<7->]\foreach \i [evaluate=\i as \j using int(\i+1)] in {0,...,5}{\draw[flecha] (j\i.east) -- (j\j.west|-j\i.east) ;}
 \draw[decorate,decoration={brace,amplitude=5pt},oro!90!ladrillo,line width=0.8pt] (-0.75,0.15) -- (6.2,0.15)
   node[midway,above=5pt,etiq]{\textbf{«Y reposó la tierra»}: cuatro veces, y ya no vuelve};\end{scope}
 \begin{scope}[visible on=<7->]\draw[sepia,dashed,-{Latex[length=2mm]}] (-0.7,-4.35) -- (11.9,-4.35) node[pos=0.5,below,etiq,text=sepia]{talla de los jueces y unidad de las tribus};\end{scope}
\end{tikzpicture}
\elabpropia
\vspace{6pt}
\begin{clave}
De Otoniel a Sansón decaen la talla de los jueces y la unidad de las tribus. Los apéndices (Jue 17--21) tocan fondo: «no había rey en Israel» (21:25).
\end{clave}
\end{frame}

\begin{frame}{El cántico de Débora (Jue 5)}
\begin{columns}[T,onlytextwidth]
\begin{column}{0.45\textwidth}
\footnotesize Débora, «profetisa» que «juzgaba a Israel» (4:4), canta una victoria en que las tribus se miden por haber acudido o no «en ayuda de YHWH» (5:23).
\begin{cita}[fontupper=\scriptsize\itshape]Los que le aman sean como el sol cuando sale en su fuerza.\fuente{(Jue 5:31)}\end{cita}
\begin{notalam}[Datación en debate]
Para muchos, junto con Éx 15, uno de los poemas más antiguos de la Biblia (Cross y Freedman, 1997); otros lo datan más tarde.
\end{notalam}
\end{column}
\begin{column}{0.51\textwidth}
\centering
\begin{tikzpicture}[chip/.style={bloque,font=\tiny,text width=1.05cm,inner sep=2.5pt}]
 \node[bloqueazul,text width=1.15cm,font=\tiny\bfseries] (A) at (0,0) {Alaba};
 \foreach \t [count=\k from 0] in {Efraín,Benjamín,Maquir,Zabulón,Isacar,Neftalí}{
   \pgfmathsetmacro{\cx}{1.6+mod(\k,3)*1.38}\pgfmathsetmacro{\cy}{0.3-floor(\k/3)*0.6}
   \node[chip] at (\cx,\cy) {\t};}
 \begin{scope}[visible on=<2->]\node[bloqueazul,text width=1.15cm,font=\tiny\bfseries,fill=oro!80!ladrillo,draw=oro!80!ladrillo] (R) at (0,-1.55) {Reprocha};\end{scope}
 \begin{scope}[visible on=<3->]\foreach \t [count=\k from 0] in {Rubén,Galaad,Dan,Aser}{
   \pgfmathsetmacro{\cx}{1.6+mod(\k,3)*1.38}\pgfmathsetmacro{\cy}{-1.25-floor(\k/3)*0.6}
   \node[chip,fill=oro!14!white,draw=oro!90!ladrillo] at (\cx,\cy) {\t};}\end{scope}
 \begin{scope}[visible on=<3->]\node[bloqueazul,text width=1.15cm,font=\tiny\bfseries,fill=oro!55!ladrillo,draw=oro!55!ladrillo] at (0,-2.6) {Maldice};\end{scope}
 \begin{scope}[visible on=<4->]\node[chip,fill=white,draw=oro!90!ladrillo,text width=2.4cm] at (2.29,-2.6) {Meroz (5:23)};\end{scope}
 \begin{scope}[visible on=<5->]\node[bloqueazul,text width=1.15cm,font=\tiny\bfseries,fill=sepia,draw=sepia] (N) at (0,-3.5) {No nombra};\end{scope}
 \begin{scope}[visible on=<5->]\node[chip,fill=white,dashed,draw=sepia] at (1.6,-3.5) {Judá};\end{scope}
 \begin{scope}[visible on=<6->]\node[chip,fill=white,dashed,draw=sepia] at (2.98,-3.5) {Simeón};\end{scope}
\end{tikzpicture}
\notafig{Elaboración propia a partir de Jue 5:14--23.}
\end{column}
\end{columns}
\end{frame}

\begin{frame}{Gedeón: «YHWH reinará»}
\begin{columns}[T,onlytextwidth]
\begin{column}{0.50\textwidth}
\centering{\heb{יְהוָה יִמְשֹׁל בָּכֶם}}\par\smallskip
\begin{cita}No reinaré yo sobre vosotros, ni mi hijo reinará sobre vosotros; YHWH reinará sobre vosotros.\fuente{(Jue 8:23)}\end{cita}
\begin{clave}
La fórmula que mejor expresa el ideal teocrático del período. La victoria sobre Madián será figura profética: «como en el día de Madián» (Is 9:4 [9:3]; 10:26).
\end{clave}
\end{column}
\begin{column}{0.46\textwidth}
\centering
\begin{tikzpicture}[node distance=0.32cm]
 \node[bloque,text width=4.6cm] (g1) {Trescientos hombres: «para que Israel no se alabe contra mí» (7:2)};
 \begin{scope}[visible on=<2->]\node[bloqueazul,text width=4.6cm,below=of g1] (g2) {Rechaza la corona (8:23)};\end{scope}
 \begin{scope}[visible on=<3->]\node[bloqueoro,text width=4.6cm,below=of g2] (g3) {Fabrica un efod: «todo Israel prostituyó tras él»; «tropiezo» (8:27)};\end{scope}
 \begin{scope}[visible on=<4->]\node[bloque,text width=4.6cm,below=of g3] (g4) {Abimelec, «mi padre es rey», se hace rey (Jue 9); fábula de Jotam};\end{scope}
 \begin{scope}[visible on=<4->]\foreach \a/\b in {g1/g2,g2/g3,g3/g4}{\draw[flecha] (\a) -- (\b);}\end{scope}
\end{tikzpicture}
\elabpropia
\end{column}
\end{columns}
\end{frame}

\begin{frame}{Jefté y Sansón}
\begin{columns}[T,onlytextwidth]
\begin{column}{0.48\textwidth}\footnotesize
\textbf{\textcolor{cafe}{Jefté (Jue 11--12)}}
\begin{itemize}
\item Conoce la historia del Éxodo (11:12--28), pero su voto sigue la lógica de los dioses vecinos.
\item \heb{שִׁבֹּלֶת}\,/\,\heb{סִבֹּלֶת}: Israel ya mata a Israel (12:6).
\end{itemize}
\begin{notalam}[Tradición judía: el voto de Jefté]
Josefo entiende que la sacrificó, y condena el acto. Kimji la ve consagrada a la virginidad (11:39).
\end{notalam}
\end{column}
\begin{column}{0.48\textwidth}\footnotesize
\textbf{\textcolor{cafe}{Sansón (Jue 13--16)}}
\begin{itemize}
\item \heb{נְזִיר אֱלֹהִים} (\emph{nezir Elohim}) desde el vientre (13:5).
\item Quebranta uno a uno los signos de su consagración: cadáver, banquete, cabello.
\item «Tómamela, porque es recta a mis ojos» (\heb{יָשְׁרָה בְעֵינָי}, 14:3).
\item Aun así, figura en Heb 11:32.
\end{itemize}
\end{column}
\end{columns}
\vspace{6pt}
\begin{clave}
En Sansón se reconoce al Israel de la época: consagrado desde su origen, vive como las naciones sin perder su llamado.
\end{clave}
\end{frame}

\begin{frame}{Un levita a sueldo (Jue 17--18)}
\begin{columns}[T,onlytextwidth]
\begin{column}{0.47\textwidth}
\begin{cita}Ahora sé que YHWH me prosperará, porque tengo un levita por sacerdote.\fuente{(Jue 17:13)}\end{cita}
\begin{notalam}[Precisión textual (18:30)]
«Jonatán, hijo de Gersón, hijo de \emph{Manasés}»: la \emph{nun} de \heb{מנשה} va suspendida; sin ella se lee \heb{משה}, «Moisés». El Talmud: se añadió para no avergonzar a Moisés (Bava Batra 109b).
\end{notalam}
\end{column}
\begin{column}{0.49\textwidth}
\centering
\begin{tikzpicture}[node distance=0.26cm]
 \node[bloque,text width=4.7cm] (l1) {Micaía: ídolo, efod y terafines (17:5)};
 \begin{scope}[visible on=<2->]\node[bloque,text width=4.7cm,below=of l1] (l2) {Un levita de Belén, por diez siclos al año};\end{scope}
 \begin{scope}[visible on=<3->]\node[bloqueoro,text width=4.7cm,below=of l2] (l3) {Dan roba el ídolo y al levita; arrasa Lais};\end{scope}
 \begin{scope}[visible on=<4->]\node[bloqueoro,text width=4.7cm,below=of l3] (l4) {Santuario en Dan «hasta el cautiverio de la tierra» (18:30)};\end{scope}
 \begin{scope}[visible on=<5->]\node[bloqueazul,text width=4.7cm,below=of l4] (l5) {Becerro de Jeroboam en Dan (1~R 12:29) $\to$ destierro asirio (Estudio~5)};\end{scope}
 \begin{scope}[visible on=<5->]\foreach \a/\b in {l1/l2,l2/l3,l3/l4,l4/l5}{\draw[flecha] (\a) -- (\b);}\end{scope}
\end{tikzpicture}
\elabpropia
\end{column}
\end{columns}
\end{frame}

\begin{frame}{No había rey en Israel}
\fotolado{panel_jue21.jpg}{Elaboración propia: Jue 21:25 en texto hebreo consonántico (BHS).}
\begin{columns}[c,onlytextwidth]
\begin{column}{0.57\textwidth}
\begin{cita}En aquellos días no había rey en Israel; cada uno hacía lo que era recto ante sus propios ojos.\fuente{(Jue 21:25; cf. 17:6)}\end{cita}
\scriptsize
\begin{itemize}
\item \textbf{Guibeá} reproduce la escena de Sodoma (Jue 19:22; Gn 19:4--5): Israel se ha vuelto Canaán.
\item \textbf{Benjamín} casi desaparece: «que falte hoy de Israel una tribu» (21:3).
\item \textbf{Dt 12:8}: \heb{אִישׁ כָּל־הַיָּשָׁר בְּעֵינָיו} frente a lo recto ante los ojos de YHWH (12:25).
\item \textbf{Os 10:9}: «desde los días de Guibeá has pecado, oh Israel».
\end{itemize}
\end{column}
\begin{column}{0.41\textwidth}\end{column}
\end{columns}
\end{frame}

\begin{frame}{Tu pueblo será mi pueblo}
\fotolado{panel_rt1.jpg}{Elaboración propia: Rt 1:16 en texto hebreo consonántico (BHS).}
\begin{columns}[c,onlytextwidth]
\begin{column}{0.57\textwidth}
\begin{cita}Adonde tú vayas, iré yo; y donde tú pases la noche, la pasaré yo. Tu pueblo será mi pueblo, y tu Dios, mi Dios.\fuente{(Rt 1:16)}\end{cita}
\scriptsize
\begin{itemize}
\item «En los días en que juzgaban los jueces» (1:1), en la Belén de Jue 17--21; aquí, en cambio, una moabita vive como israelita fiel.
\item Un orden significativo: \destacar{primero el pueblo, después Dios}.
\item \heb{חֶסֶד} (\emph{jésed}): lealtad más allá de lo debido (1:8; 2:20; 3:10).
\item \heb{כָּנָף}: «bajo sus alas» (2:12) y «extiende tu manto» (3:9).
\end{itemize}
\end{column}
\begin{column}{0.41\textwidth}\end{column}
\end{columns}
\end{frame}

\begin{frame}{Booz, el \emph{go'el}}
\begin{columns}[T,onlytextwidth]
\begin{column}{0.52\textwidth}\footnotesize
\begin{itemize}
\item \heb{גֹּאֵל} (\emph{go'el}): rescata la tierra y el nombre del pariente (Lv 25:25; cf. Dt 25:5--10).
\item «No te ha dejado hoy sin redentor» (Rt 4:14).
\item \heb{הוּא אֲבִי־יִשַׁי אֲבִי דָוִד}: «padre de Isaí, padre de David» (4:17).
\item Rut en la genealogía del Mesías (Mt 1:5).
\end{itemize}
\vspace{4pt}
\begin{clave}
La genealogía convierte una historia familiar en prólogo de la monarquía y, para Mateo, del Mesías.
\end{clave}
\end{column}
\begin{column}{0.44\textwidth}
\centering
\begin{tikzpicture}[g/.style={bloque,text width=1.5cm,inner sep=2.5pt},node distance=0.32cm]
 \node[g] (f) {Fares};
 \begin{scope}[visible on=<2->]\node[g,below=0.75cm of f] (b) {Booz};\end{scope}
 \begin{scope}[visible on=<3->]\node[g,right=0.45cm of b,fill=oro!14!white,draw=oro!90!ladrillo] (r) {Rut, la moabita};\end{scope}
 \begin{scope}[visible on=<3->]\draw[oro!90!ladrillo,line width=1pt] (b) -- (r);\end{scope}
 \begin{scope}[visible on=<3->]\draw[flecha,dashed] (f) -- node[right,etiq,text=sepia]{\ldots} (b);\end{scope}
 \begin{scope}[visible on=<4->]\node[g,below=0.5cm of b] (o) {Obed};\end{scope}
 \begin{scope}[visible on=<5->]\node[g,below=of o] (i) {Isaí};\end{scope}
 \begin{scope}[visible on=<6->]\node[bloqueazul,text width=1.5cm,below=of i] (d) {David};\end{scope}
 \begin{scope}[visible on=<6->]\draw[flecha] (b.south) -- (o.north);\end{scope} \begin{scope}[visible on=<6->]\draw[flecha] (o) -- (i);\end{scope} \begin{scope}[visible on=<6->]\draw[flecha] (i) -- (d);\end{scope}
 \begin{scope}[visible on=<6->]\draw[decorate,decoration={brace,amplitude=4pt},sepia] ([xshift=0.2cm]r.east|-f.north) -- ([xshift=0.2cm]r.east|-d.south)
   node[midway,right=5pt,etiq,align=left]{diez\\generaciones\\(Rt 4:18--22)};\end{scope}
\end{tikzpicture}
\notafig{Elaboración propia a partir de Rt 4:17 y Mt 1:5.}
\end{column}
\end{columns}
\end{frame}

\begin{frame}{¿Y Dt 23:3 [23:4]?}
\begin{columns}[T,onlytextwidth]
\begin{column}{0.47\textwidth}
\begin{cita}No entrará amonita ni moabita en la congregación de YHWH.\fuente{(Dt 23:3 [23:4])}\end{cita}
\begin{notalam}[Tradición judía: «moabita, y no moabitesa»]
La Mishná limita la prohibición a los varones (Yevamot 8:3) y el Talmud lo condensa en esa frase (Yevamot 76b--77a), con la que zanja la duda de Doeg el edomita sobre la legitimidad de David. Es interpretación rabínica, no distinción explícita de la Torá.
\end{notalam}
\end{column}
\begin{column}{0.49\textwidth}
\centering
\begin{tikzpicture}
 \node[bloqueazul,text width=5.1cm] (t) at (0,0) {Dt 23:3 [23:4]\\[-1pt]{\mdseries\scriptsize «no entrará amonita ni moabita»}};
 \begin{scope}[visible on=<2->]\node[bloqueoro,text width=2.3cm,font=\scriptsize,anchor=north] (m) at (-1.4,-0.9) {\textbf{Lectura rabínica}\\solo varones\\(Yevamot 8:3; 76b--77a)};\end{scope}
 \begin{scope}[visible on=<3->]\node[bloque,text width=2.3cm,font=\scriptsize,anchor=north] (e) at (1.4,-0.9) {\textbf{Lectura rigurosa}\\aplicación estricta\\(Esd 9--10; Neh 13:1--3)};\end{scope}
 \begin{scope}[visible on=<3->]\draw[flecha] (t.south) -- (m.north);\end{scope} \begin{scope}[visible on=<3->]\draw[flecha] (t.south) -- (e.north);\end{scope}
 \begin{scope}[visible on=<4->]\node[bloque,fill=white,dashed,text width=5.1cm,anchor=north] (r) at (0,-2.75) {Rut, la moabita, en el linaje de David (Rt 4:17)};\end{scope}
 \begin{scope}[visible on=<4->]\draw[flecha,dashed] (m.south) -- (m.south|-r.north);\end{scope}
\end{tikzpicture}
\elabpropia
\end{column}
\end{columns}
\end{frame}

% ===================================================
\section{Samuel}
\begin{frame}{El cántico de Ana}
\fotolado{panel_1s2.jpg}{Elaboración propia: 1~S 2:10 en texto hebreo consonántico (BHS).}
\begin{columns}[c,onlytextwidth]
\begin{column}{0.57\textwidth}
\begin{cita}YHWH juzgará los confines de la tierra; dará poder a su rey y exaltará el poder de su ungido.\fuente{(1~S 2:10)}\end{cita}
\scriptsize
\begin{itemize}
\item Una estéril, como Sara: Dios la «visita» (\heb{פָּקַד}, 1~S 2:21; Gn 21:1).
\item Poema de la inversión: «levanta del polvo al pobre» (2:8).
\item Habla de «su rey» cuando Israel aún no tiene ninguno.
\item \heb{מְשִׁיחוֹ} (\emph{meshijó}): primera vez que el «ungido» es un rey.
\end{itemize}
\end{column}
\begin{column}{0.41\textwidth}\end{column}
\end{columns}
\end{frame}

\begin{frame}{«Su ungido»: del hebreo a Lucas}
\centering
\begin{tikzpicture}[node distance=0.9cm]
\node[caja,text width=3.3cm,align=center] (h) {\textbf{Hebreo (1~S 2:10)}\\[3pt]\heb{וְיָרֵם קֶרֶן מְשִׁיחוֹ}\\[2pt]«exaltará el cuerno\\de su ungido»};
\begin{scope}[visible on=<2->]\node[caja,text width=3.3cm,align=center,right=of h] (g) {\textbf{LXX}\\[3pt]\grc{ὑψώσει κέρας χριστοῦ αὐτοῦ}\\[2pt]\emph{mashíaj} $\to$ \grc{χριστός}};\end{scope}
\begin{scope}[visible on=<3->]\node[caja,text width=3.3cm,align=center,right=of g] (p) {\textbf{Lucas (1:69)}\\[3pt]\grc{ἤγειρεν κέρας σωτηρίας}\\[2pt]«en la casa de David»};\end{scope}
\begin{scope}[visible on=<3->]\draw[-{Latex[length=2.5mm]},oro!90!ladrillo,line width=1.4pt] (h) -- (g);\end{scope}
\begin{scope}[visible on=<3->]\draw[-{Latex[length=2.5mm]},oro!90!ladrillo,line width=1.4pt] (g) -- (p);\end{scope}
\end{tikzpicture}
\elabpropia
\vspace{4pt}
{\small\itshape\color{sepia}\centering María canta como Ana: «quitó de los tronos a los poderosos y exaltó a los humildes» (Lc 1:52; cf. 1~S 2:7--8).\par}
\vspace{4pt}
\begin{notalam}[Crítica textual]
4QSam\textsuperscript{a} añade en 1~S 1:22 que Samuel será «nazareo para siempre»; la LXX de 1~S 2:10 inserta un texto cercano a Jer 9:23--24 [9:22--23]. Aquí se traduce el texto masorético.
\end{notalam}
\end{frame}

\begin{frame}{Icabod y Eben-ezer}
\begin{columns}[T,onlytextwidth]
\begin{column}{0.48\textwidth}
\centering{\heb{אִי־כָבוֹד \foreignlanguage{spanish}{[\ldots]} גָּלָה כָבוֹד מִיִּשְׂרָאֵל}}\par\smallskip
\begin{cita}Y llamó al niño Icabod, diciendo: ¡Se ha ido la gloria de Israel!\fuente{(1~S 4:21)}\end{cita}
\begin{notalam}[Filología: \heb{גָּלָה}]
El verbo del destierro: la gloria «va al exilio». Los hijos de Elí «no conocían a YHWH» (2:12), y el sacerdocio arrastra al pueblo.
\end{notalam}
\end{column}
\begin{column}{0.48\textwidth}
\centering{\heb{עַד־הֵנָּה עֲזָרָנוּ יְהוָה}}\par\smallskip
\begin{cita}Y le puso por nombre Eben-ezer, diciendo: Hasta aquí nos ayudó YHWH.\fuente{(1~S 7:12)}\end{cita}
\begin{clave}
\textbf{Una promesa en la ruina.} 1~S 2:35 anuncia «un sacerdote fiel» que andará «delante de mi ungido».
\end{clave}
\end{column}
\end{columns}
\end{frame}

\begin{frame}{Afec y Mizpa: dos maneras de acudir a Dios}
\centering
\begin{tikzpicture}[x=1cm,y=1cm]
 \node[bloqueazul,text width=1.7cm] (A) at (0,0) {Afec\\[-1pt]{\tiny\mdseries 1~S 4--6}};
 \node[bloque,text width=2.9cm] (a1) at (2.75,0) {El arca como amuleto: «para que nos salve» (4:3)};
 \begin{scope}[visible on=<2->]\node[bloque,text width=2.6cm] (a2) at (6.0,0) {Derrota; los filisteos capturan el arca};\end{scope}
 \begin{scope}[visible on=<3->]\node[bloque,text width=2.7cm,fill=white] (a3) at (9.2,0) {\textbf{Icabod}: «se ha ido la gloria» (4:21)};\end{scope}
 \begin{scope}[visible on=<4->]\node[bloqueazul,text width=1.7cm] (M) at (0,-1.9) {Mizpa\\[-1pt]{\tiny\mdseries 1~S 7}};\end{scope}
 \begin{scope}[visible on=<4->]\node[bloqueoro,text width=2.9cm] (m1) at (2.75,-1.9) {«Quitad los dioses ajenos» (7:3)};\end{scope}
 \begin{scope}[visible on=<5->]\node[bloqueoro,text width=2.6cm] (m2) at (6.0,-1.9) {Ayuno, confesión e intercesión de Samuel};\end{scope}
 \begin{scope}[visible on=<6->]\node[bloqueoro,text width=2.7cm,fill=white] (m3) at (9.2,-1.9) {\textbf{Eben-ezer}: «hasta aquí nos ayudó YHWH» (7:12)};\end{scope}
 \begin{scope}[visible on=<6->]\foreach \a/\b in {A/a1,a1/a2,a2/a3,M/m1,m1/m2,m2/m3}{\draw[flecha] (\a) -- (\b);}\end{scope}
 \begin{scope}[visible on=<6->]\draw[dashed,sepia] (-1,-0.95) -- (10.6,-0.95);\end{scope}
\end{tikzpicture}
\elabpropia
\vspace{8pt}
\begin{clave}
Con el arca como objeto, Israel perdió; al volver a Dios, fue librado. Dagón cae ante el arca, y el arca vuelve sola a Bet-semes (1~S 5--6): la gloria no depende del objeto.
\end{clave}
\end{frame}

\begin{frame}{«Como todas las naciones»}
\centering
\begin{tikzpicture}
 \foreach \x/\n/\t/\d/\r in {0/1/{Éx 19:6}/{«Reino de sacerdotes y nación santa»: un pueblo distinto para las naciones.}/{Nm 23:9},
   4.2/2/{Dt 17:14--20}/{Un rey \emph{que Dios escoja}, sin caballos, mujeres ni oro, que lea la Torá cada día.}/{Gn 17:6; 49:10},
   8.4/3/{1~S 8:5, 20}/{\heb{כְּכָל־הַגּוֹיִם}: «un rey [\ldots] como todas las naciones».}/{8:7: «a mí me han desechado»}}{
  \node[caja,anchor=north] (b\n) at (\x,0) {\parbox[t][2.75cm][t]{3.55cm}{\RaggedRight
    {\small\textbf{\textcolor{ladrillo}{\t}}}\\[5pt]\small \d\par\vfill{\scriptsize\color{sepia}(\r)}}};
  \node[circle,fill=cafe,text=piedra,font=\scriptsize\bfseries,inner sep=1.6pt] at (b\n.north west) {\n};}
\end{tikzpicture}
\elabpropia
\vspace{8pt}
\begin{clave}
\small La Torá prevé la monarquía; lo que está en juego es el \destacar{modelo}: un rey a la manera de las naciones, en lugar de un rey sujeto a la Torá en un pueblo llamado a ser distinto.
\end{clave}
\end{frame}

% ===================================================
\section{Saúl y David}
\begin{frame}{Saúl: obedecer es mejor}
\begin{columns}[T,onlytextwidth]
\begin{column}{0.44\textwidth}
\centering
\begin{tikzpicture}[node distance=0.24cm]
 \node[bloque,text width=4.3cm] (s1) {Ungido \heb{נָגִיד} (\emph{naguid}) sobre la heredad de YHWH (10:1)};
 \begin{scope}[visible on=<2->]\node[bloque,text width=4.3cm,below=of s1] (s2) {Libra a Jabes, la ciudad arrasada en Jue 21 (1~S 11)};\end{scope}
 \begin{scope}[visible on=<3->]\node[bloqueoro,text width=4.3cm,below=of s2] (s3) {Gilgal: ofrece él mismo el holocausto (13:8--12)};\end{scope}
 \begin{scope}[visible on=<4->]\node[bloqueoro,text width=4.3cm,below=of s3] (s4) {Amalec: perdona a Agag (1~S 15)};\end{scope}
 \begin{scope}[visible on=<5->]\node[bloqueazul,text width=4.3cm,below=of s4] (s5) {El reino es «rasgado» (15:28)};\end{scope}
 \begin{scope}[visible on=<5->]\foreach \a/\b in {s1/s2,s2/s3,s3/s4,s4/s5}{\draw[flecha] (\a) -- (\b);}\end{scope}
\end{tikzpicture}
\elabpropia
\end{column}
\begin{column}{0.52\textwidth}
\centering{\heb{הִנֵּה שְׁמֹעַ מִזֶּבַח טוֹב}}\par\smallskip
\begin{cita}Ciertamente obedecer es mejor que el sacrificio, y prestar atención, mejor que la grasa de los carneros.\fuente{(1~S 15:22)}\end{cita}
\begin{cita}[fontupper=\scriptsize\itshape]YHWH no desamparará a su pueblo, por causa de su gran nombre.\fuente{(1~S 12:22)}\end{cita}
\begin{notalam}[Griego: de la LXX a Romanos]
\grc{οὐκ ἀπώσεται κύριος τὸν λαὸν αὐτοῦ}: Pablo usa el mismo verbo en Ro 11:1--2. La permanencia descansa en el nombre de Dios, no en la fidelidad del pueblo.
\end{notalam}
\end{column}
\end{columns}
\end{frame}

\begin{frame}{Perla filológica · «no desechará a su pueblo»}
{\footnotesize La promesa de Samuel, el día en que el pueblo pide rey, llega por la Septuaginta al centro de Romanos 11.}\par\vspace{4pt}
\centering\scriptsize\renewcommand{\arraystretch}{1.22}
\begin{tabular}{>{\RaggedRight\arraybackslash}p{2.1cm}>{\RaggedRight\arraybackslash}p{5.0cm}>{\RaggedRight\arraybackslash}p{4.3cm}}
\toprule
{\color{ladrillo}\bfseries Texto} & {\color{ladrillo}\bfseries Forma} & {\color{ladrillo}\bfseries Observación}\\\midrule
1~S 12:22 TM & \heb{כִּי לֹא־יִטֹּשׁ יְהוָה אֶת־עַמּוֹ בַּעֲבוּר שְׁמוֹ הַגָּדוֹל} & La razón: el nombre de Dios\\
1~S 12:22 LXX & \grc{ὅτι οὐκ ἀπώσεται κύριος τὸν λαὸν αὐτοῦ} & «Abandonar» se vuelve «rechazar» (\grc{ἀπωθέω})\\
Sal 94:14 [93:14] & \grc{ὅτι οὐκ ἀπώσεται κύριος τὸν λαὸν αὐτοῦ} & La fórmula, ya confesión del salterio\\
Ro 11:1--2 & \grc{μὴ ἀπώσατο ὁ θεὸς τὸν λαὸν αὐτοῦ;} \ldots{} \grc{οὐκ ἀπώσατο} & Pablo pregunta y responde con la promesa\\
\bottomrule
\end{tabular}
\elabpropia
\vspace{4pt}
\begin{clave}\footnotesize La respuesta de Pablo, «de ninguna manera», repite lo que Samuel dijo a un pueblo que acababa de rechazar a Dios como rey.\end{clave}
\end{frame}

\begin{frame}{David: el menor de Belén}
\begin{columns}[T,onlytextwidth]
\begin{column}{0.48\textwidth}
\centering{\heb{וַיהוָה יִרְאֶה לַלֵּבָב}}\par\smallskip
\begin{cita}El hombre mira lo que está delante de sus ojos, pero YHWH mira el corazón.\fuente{(1~S 16:7)}\end{cita}
\footnotesize
\begin{itemize}
\item Contrapunto al estribillo de Jueces: Dios elige mirando el corazón.
\item El menor, con las ovejas; el Espíritu viene sobre él «desde aquel día» (16:13).
\item El ungido de 2:10 recibe un nombre.
\end{itemize}
\end{column}
\begin{column}{0.48\textwidth}
\centering{\heb{כִּי יֵשׁ אֱלֹהִים לְיִשְׂרָאֵל}}\par\smallskip
\begin{cita}Toda la tierra sabrá que hay Dios en Israel [\ldots], porque de YHWH es la batalla.\fuente{(1~S 17:46--47)}\end{cita}
\begin{notalam}[Crítica textual: la estatura de Goliat]
«Seis codos y un palmo» en el TM; «cuatro codos y un palmo» en LXX, 4QSam\textsuperscript{a} y Josefo. Cf. 2~S 21:19 y 1~Cr 20:5.
\end{notalam}
\end{column}
\end{columns}
\end{frame}

\begin{frame}{Los ojos del hombre y los ojos de YHWH}
\centering
\begin{tikzpicture}[x=1cm,y=1cm]
 \node[bloqueazul,text width=3.3cm] (h) at (-3.0,0) {«Lo recto ante sus ojos»\\[-1pt]{\tiny\mdseries\heb{הַיָּשָׁר בְּעֵינָיו}}};
 \node[bloqueazul,text width=3.3cm,fill=oro!80!ladrillo,draw=oro!80!ladrillo] (y) at (3.0,0) {«YHWH mira el corazón»\\[-1pt]{\tiny\mdseries\heb{וַיהוָה יִרְאֶה לַלֵּבָב}}};
 \begin{scope}[visible on=<2->]\node[bloque,text width=3.6cm,font=\scriptsize,below=0.45cm of h] (h1) {Sansón: «recta a mis ojos» (Jue 14:3)};\end{scope}
 \begin{scope}[visible on=<3->]\node[bloque,text width=3.6cm,font=\scriptsize,below=0.2cm of h1] (h2) {«Cada uno hacía lo recto ante sus ojos» (Jue 17:6; 21:25)};\end{scope}
 \begin{scope}[visible on=<4->]\node[bloque,text width=3.6cm,font=\scriptsize,below=0.2cm of h2] (h3) {Saúl se justifica ante Samuel (1~S 15:20--21)};\end{scope}
 \begin{scope}[visible on=<4->]\node[bloqueoro,text width=3.6cm,font=\scriptsize,below=0.45cm of y] (y1) {Eliab, el primogénito, descartado (1~S 16:6--7)};\end{scope}
 \begin{scope}[visible on=<5->]\node[bloqueoro,text width=3.6cm,font=\scriptsize,below=0.2cm of y1] (y2) {El menor, con las ovejas, ungido (16:11--13)};\end{scope}
 \begin{scope}[visible on=<6->]\node[bloqueoro,text width=3.6cm,font=\scriptsize,below=0.2cm of y2] (y3) {El Espíritu «desde aquel día» (16:13)};\end{scope}
 \begin{scope}[visible on=<6->]\draw[flecha,line width=1.4pt] (h.east) -- node[above,etiq]{Dt 12:8 / 12:25} (y.west);\end{scope}
 \begin{scope}[visible on=<6->]\draw[sepia,dashed] (0,-0.75) -- (0,-3.5);\end{scope}
\end{tikzpicture}
\notafig{Elaboración propia a partir de Dt 12:8, 25; Jue 17:6; 21:25; 1~S 16:6--13.}
\vspace{4pt}
\begin{clave}
A un pueblo que hacía «lo recto ante sus ojos», Dios le opone una elección basada en lo que ve en el corazón.
\end{clave}
\end{frame}

\begin{frame}{Saúl y David frente a frente}
\centering\footnotesize\renewcommand{\arraystretch}{1.25}
\begin{tabular}{>{\RaggedRight\arraybackslash}p{2.4cm}>{\RaggedRight\arraybackslash}p{4.6cm}>{\RaggedRight\arraybackslash}p{4.6cm}}
\toprule
& {\color{ladrillo}\bfseries Saúl} & {\color{ladrillo}\bfseries David}\\\midrule
Tribu & Benjamín (9:1--2) & Judá, Belén (16:1)\\
Elección & El pueblo pide un rey «como las naciones» (8:5) & «Un hombre conforme a su corazón» (13:14)\\
Ante el sacerdocio & Mata a los sacerdotes de Nob (22:18) & Respeta al «ungido de YHWH» (24:6 [24:7])\\
Ante el pecado & Se justifica (15:20) & Confiesa (2~S 12:13)\\
Final & Endor y Gilboa (1~S 28; 31) & Rey de la casa de Judá en Hebrón (2~S 2:4)\\
\bottomrule
\end{tabular}
\notafig{Elaboración propia a partir de 1--2 Samuel.}
\vspace{6pt}
\begin{notalam}[Tradición judía: Yoma 22b]
«Saúl falló en una cosa, y se le contó; David en dos, y no se le contó». El propio texto insinúa la razón: Saúl se justifica (1~S 15:20--21); David confiesa (2~S 12:13).
\end{notalam}
\end{frame}

\begin{frame}{Jonatán, Gilboa y dos casas}
\begin{columns}[T,onlytextwidth]
\begin{column}{0.48\textwidth}\footnotesize
\begin{itemize}
\item \textbf{Jonatán}: «tú reinarás sobre Israel, y yo seré segundo» (23:17).
\item \textbf{Nob}: Saúl mata a los sacerdotes (22:18; cf. 2:31--33).
\item \textbf{David perdona} al «ungido de YHWH» (\heb{מְשִׁיחַ יְהוָה}, 24:6 [24:7]).
\item \textbf{Abigail}: Dios le hará «casa firme» (\heb{בַּיִת נֶאֱמָן}, 25:28).
\item \textbf{Endor y Gilboa} (1~S 28; 31): «¡Cómo han caído los valientes!» (2~S 1:19).
\end{itemize}
\end{column}
\begin{column}{0.48\textwidth}
\centering
\begin{tikzpicture}
 \node[bloque,text width=4.4cm] (g) at (0,0) {Muerte de Saúl en Gilboa (1~S 31)};
 \begin{scope}[visible on=<2->]\node[bloqueazul,text width=1.95cm,font=\tiny\bfseries,anchor=north] (j) at (-1.22,-0.75) {Casa de Judá\\[1pt]{\mdseries David, en Hebrón}\\[1pt]{\mdseries (2~S 2:4)}};\end{scope}
 \begin{scope}[visible on=<3->]\node[bloqueazul,text width=1.95cm,font=\tiny\bfseries,anchor=north,fill=sepia,draw=sepia] (i) at (1.22,-0.75) {«Todo Israel»\\[1pt]{\mdseries Is-boset}\\[1pt]{\mdseries (2~S 2:9--10)}};\end{scope}
 \begin{scope}[visible on=<3->]\draw[flecha] (g.south) -- (j.north);\end{scope} \begin{scope}[visible on=<3->]\draw[flecha] (g.south) -- (i.north);\end{scope}
 \begin{scope}[visible on=<4->]\node[bloqueoro,text width=4.4cm,anchor=north] (u) at (0,-2.35) {David reúne ambas casas (2~S 5)\\{\tiny división definitiva: Estudio~4}};\end{scope}
 \begin{scope}[visible on=<4->]\draw[flecha] (j.south) -- (j.south|-u.north);\end{scope} \begin{scope}[visible on=<4->]\draw[flecha] (i.south) -- (i.south|-u.north);\end{scope}
\end{tikzpicture}
\elabpropia
\end{column}
\end{columns}
\vspace{4pt}
\begin{clave}
\textbf{Judá e Israel.} Por primera vez, dos entidades políticas, ya insinuadas en 1~S 11:8 y 18:16. Aquí «Israel» es el conjunto de tribus fieles a la casa de Saúl, no todavía el reino del norte.
\end{clave}
\end{frame}

% ===================================================
\section{Vocación}
\begin{frame}{La vocación quebrada y preservada}
\centering
\begin{tikzpicture}[x=1cm,y=1cm]
 \node[bloqueazul,text width=11.2cm] (top) at (0,0) {Éx 19:6: «reino de sacerdotes y nación santa»};
 \node[font=\scriptsize\bfseries,text=ladrillo] (Q) at (-3.05,-0.65) {Quebrada};
 \node[font=\scriptsize\bfseries,text=oro!70!ladrillo] (P) at (3.05,-0.65) {Preservada: el remanente};
 \begin{scope}[visible on=<2->]\node[bloque,text width=5.6cm,font=\scriptsize,inner sep=2.5pt,below=0.12cm of Q] (q1) {\textbf{Cananeización}: los baales (Jue 2:11--13; cf. Nm 22--25)};\end{scope}
 \begin{scope}[visible on=<2->]\node[bloque,text width=5.6cm,font=\scriptsize,inner sep=2.5pt,below=0.18cm of q1] (q2) {\textbf{Sacerdocio vendido}: el levita de Micaía; los hijos de Elí};\end{scope}
 \begin{scope}[visible on=<3->]\node[bloque,text width=5.6cm,font=\scriptsize,inner sep=2.5pt,below=0.18cm of q2] (q3) {\textbf{Culto como magia}: el arca amuleto; Gilgal; Endor};\end{scope}
 \begin{scope}[visible on=<3->]\node[bloque,text width=5.6cm,font=\scriptsize,inner sep=2.5pt,below=0.18cm of q3] (q4) {\textbf{Israel como Sodoma}: Guibeá};\end{scope}
 \begin{scope}[visible on=<4->]\node[bloque,text width=5.6cm,font=\scriptsize,inner sep=2.5pt,below=0.18cm of q4] (q5) {\textbf{Como las naciones}: 1~S 8 frente a Nm 23:9};\end{scope}
 \begin{scope}[visible on=<4->]\node[bloqueoro,text width=5.6cm,font=\scriptsize,inner sep=2.5pt,below=0.12cm of P] (p1) {\textbf{Débora}: canta la victoria como obra de YHWH};\end{scope}
 \begin{scope}[visible on=<5->]\node[bloqueoro,text width=5.6cm,font=\scriptsize,inner sep=2.5pt,below=0.18cm of p1] (p2) {\textbf{Rut}: la bendición de Abraham a las familias (Gn 12:3)};\end{scope}
 \begin{scope}[visible on=<5->]\node[bloqueoro,text width=5.6cm,font=\scriptsize,inner sep=2.5pt,below=0.18cm of p2] (p3) {\textbf{Ana} profetiza al ungido; \textbf{Samuel} intercede (12:23)};\end{scope}
 \begin{scope}[visible on=<6->]\node[bloqueoro,text width=5.6cm,font=\scriptsize,inner sep=2.5pt,below=0.18cm of p3] (p4) {\textbf{Jonatán}: acepta la elección de Dios};\end{scope}
 \begin{scope}[visible on=<6->]\node[bloqueoro,text width=5.6cm,font=\scriptsize,inner sep=2.5pt,below=0.18cm of p4] (p5) {\textbf{David}: «toda la tierra sabrá» (17:46)};\end{scope}
 \begin{scope}[visible on=<6->]\draw[flecha] (top.south-|Q) -- (Q.north);\end{scope}
 \begin{scope}[visible on=<6->]\draw[flecha] (top.south-|P) -- (P.north);\end{scope}
\end{tikzpicture}
\elabpropia
\vspace{6pt}
\begin{clave}
«Aun en este tiempo ha quedado un remanente escogido por gracia» (Ro 11:5).
\end{clave}
\end{frame}

% ===================================================
\section{Tradición judía}
\begin{frame}{Rut y Ana en la tradición}
\small Elaboración rabínica, que conviene distinguir del texto bíblico.\par\medskip
\centering
\begin{tikzpicture}
 \foreach \x/\n/\t/\d/\r in {0/1/{Moabitesa}/{«Moabita, y no moabitesa»: la ley de Dt 23:3 [23:4] afecta a los varones.}/{Yevamot 76b--77a},
   3.15/2/{El rollo del \emph{jésed}}/{Rut no trae leyes: enseña «la recompensa de los que practican la misericordia».}/{Rut Rabbah 2:14},
   6.3/3/{Rut en Shavuot}/{Siega, aceptación de la Torá; David nació y murió en Shavuot.}/{T. de Jerusalén, Jagigá 2:3; Soferim 14},
   9.45/4/{Ana}/{Modelo de oración: corazón dirigido, labios que se mueven, voz baja. Primera en decir «YHWH de los ejércitos».}/{Berajot 31a--b}}{
  \node[caja,anchor=north] (b\n) at (\x,0) {\parbox[t][3.25cm][t]{2.65cm}{\RaggedRight
    {\small\textbf{\textcolor{ladrillo}{\t}}}\\[4pt]\scriptsize \d\par\vfill{\tiny\color{sepia}(\r)}}};
  \node[circle,fill=cafe,text=piedra,font=\scriptsize\bfseries,inner sep=1.6pt] at (b\n.north west) {\n};}
\end{tikzpicture}
\notafig{Elaboración propia a partir de las fuentes rabínicas citadas.}
\end{frame}

\begin{frame}{Jueces, Saúl y David en la tradición}
\begin{columns}[T,onlytextwidth]
\begin{column}{0.48\textwidth}\footnotesize
\begin{itemize}
\item \textbf{Sansón}: «fue tras sus ojos; por eso le sacaron los ojos» (Mishná, Sota 1:8).
\item \textbf{Cada generación}: «Jefté en su generación es como Samuel en la suya» (Rosh Hashaná 25b).
\item \textbf{Orfa}: antepasada de Goliat (Sota 42b); \textbf{Booz} = Ibzán (Bava Batra 91a).
\item \textbf{Profetisas}: Débora, Ana y Abigail (Meguilá 14a).
\end{itemize}
\end{column}
\begin{column}{0.48\textwidth}
\begin{notalam}[Autoría según el Talmud]
«Samuel escribió su libro, Jueces y Rut» (Bava Batra 14b); como 1~Samuel narra su muerte (1~S 25:1), lo completaron «Gad el vidente y Natán el profeta» (15a). La crítica moderna sitúa la redacción final mucho después.
\end{notalam}
\footnotesize
\begin{itemize}
\item \textbf{Targum de Rut}: diez hambres (1:1).
\item \textbf{Josefo}, \emph{Antigüedades} 5--6: elogia el valor de Saúl en Gilboa.
\end{itemize}
\end{column}
\end{columns}
\end{frame}

\begin{frame}{En la sinagoga y en la oración}
\begin{columns}[T,onlytextwidth]
\begin{column}{0.60\textwidth}
\scriptsize\renewcommand{\arraystretch}{1.25}
\begin{tabular}{>{\RaggedRight\arraybackslash}p{1.95cm}>{\RaggedRight\arraybackslash}p{1.7cm}>{\RaggedRight\arraybackslash}p{2.55cm}}
\toprule
{\color{ladrillo}\bfseries Ocasión} & {\color{ladrillo}\bfseries Haftará} & {\color{ladrillo}\bfseries Conexión}\\\midrule
Beshalaj & Jue 4:4--5:31 & Débora y Éx 15\\
Nasó & Jue 13:2--25 & Sansón y el nazareo (Nm 6)\\
Jukat & Jue 11:1--33 & Jefté y Nm 21\\
Kóraj & 1~S 11:14--12:22 & Samuel y Moisés; «no desamparará»\\
Rosh Hashaná 1 & 1~S 1:1--2:10 & Ana y Sara\\
Víspera de luna nueva & 1~S 20:18--42 & Jonatán y David\\
Shavuot & Rut & Siega y Torá\\
\bottomrule
\end{tabular}
\elabpropia
\end{column}
\begin{column}{0.36\textwidth}
\centering{\heb{אֶת־צֶמַח דָּוִד עַבְדְּךָ}}\par\smallskip
\begin{cita}[fontupper=\scriptsize\itshape]Haz brotar pronto el renuevo de David tu siervo.\fuente{(Amidá, 15.ª bendición)}\end{cita}
\begin{notalam}[Liturgia]
Con el lenguaje de Jer 23:5. La bendición después de comer pide por «el reino de la casa de David, tu ungido» (Sacks, 2009).
\end{notalam}
\end{column}
\end{columns}
\end{frame}

% ===================================================
\section{Restauración}
\begin{frame}{La promesa en esta etapa}
\centering\scriptsize\renewcommand{\arraystretch}{1.25}
\begin{tabular}{>{\RaggedRight\arraybackslash}p{2.0cm}>{\RaggedRight\arraybackslash}p{4.4cm}>{\RaggedRight\arraybackslash}p{5.0cm}}
\toprule
{\color{ladrillo}\bfseries Texto} & {\color{ladrillo}\bfseries Promesa} & {\color{ladrillo}\bfseries Proyección}\\\midrule
Jue 2:16, 18 & Libertadores por compasión & Is 1:26\\
Jue 7--8 & La victoria sobre Madián & Is 9:4 [9:3]; 10:26\\
Rt 4:14--17 & Un redentor y el linaje de David & Is 59:20; Ro 11:26\\
1~S 2:10 & El poder de su ungido & Sal 2:2; Lc 1:69\\
1~S 2:35 & Sacerdote fiel, casa firme & 1~S 25:28; 2~S 7:16 (Est. 4)\\
1~S 12:22 & No desamparará a su pueblo & Sal 94:14; Ro 11:1--2\\
1~S 13:14; 16:7 & Un hombre según su corazón & Jer 3:15; Hch 13:22\\
1~S 16:13 & El ungido de Belén, hijo de Isaí & Is 11:1, 10; Mi 5:2 [5:1]; Ro 15:12\\
1~S 17:46 & «Toda la tierra sabrá» & Is 55:3--5\\
\bottomrule
\end{tabular}
\elabpropia
\vspace{2pt}
{\small\itshape\color{sepia}\centering YHWH «se compadecía» de su gemido (Jue 2:18): una compasión que no se apoya en el mérito.\par}
\end{frame}

\begin{frame}{Una cadena de promesas}
\centering
\begin{tikzpicture}[x=1cm,y=1cm,col/.style={font=\scriptsize\bfseries,text=ladrillo}]
 \node[col] at (0,0.55) {En la etapa};
 \node[col] at (4.4,0.55) {Los profetas};
 \node[col] at (8.8,0.55) {Pablo};
 \begin{scope}[visible on=<2->]\node[bloque,text width=3.0cm] (a1) at (0,0) {Jueces levantados por compasión (Jue 2:18)};\end{scope}
 \begin{scope}[visible on=<2->]\node[bloque,text width=3.0cm] (a2) at (0,-1.3) {«No desamparará a su pueblo» (1~S 12:22)};\end{scope}
 \begin{scope}[visible on=<3->]\node[bloque,text width=3.0cm] (a3) at (0,-2.6) {El ungido de Belén, hijo de Isaí (1~S 16:13)};\end{scope}
 \begin{scope}[visible on=<3->]\node[bloque,text width=3.0cm] (a4) at (0,-3.9) {Booz, el \emph{go'el} (Rt 4:14)};\end{scope}
 \begin{scope}[visible on=<4->]\node[bloqueoro,text width=3.0cm] (b1) at (4.4,0) {«Restauraré tus jueces» (Is 1:26)};\end{scope}
 \begin{scope}[visible on=<4->]\node[bloqueoro,text width=3.0cm] (b2) at (4.4,-1.3) {Eco en Sal 94:14};\end{scope}
 \begin{scope}[visible on=<4->]\node[bloqueoro,text width=3.0cm] (b3) at (4.4,-2.6) {Vara de Isaí; Renuevo; «David su rey»\\{\tiny(Is 11:1; Jer 23:5; Os 3:5)}};\end{scope}
 \begin{scope}[visible on=<5->]\node[bloqueoro,text width=3.0cm] (b4) at (4.4,-3.9) {«Vendrá de Sion el Libertador»: \emph{go'el} (Is 59:20)};\end{scope}
 \begin{scope}[visible on=<5->]\node[bloqueazul,text width=3.0cm] (c2) at (8.8,-1.3) {«¿Ha desechado Dios a su pueblo?» (Ro 11:1--2)};\end{scope}
 \begin{scope}[visible on=<6->]\node[bloqueazul,text width=3.0cm] (c3) at (8.8,-2.6) {Raíz de Isaí (Ro 15:12)};\end{scope}
 \begin{scope}[visible on=<6->]\node[bloqueazul,text width=3.0cm] (c4) at (8.8,-3.9) {«Todo Israel será salvo» (Ro 11:26)};\end{scope}
 \begin{scope}[visible on=<6->]\foreach \i in {1,2,3,4}{\draw[flecha] (a\i) -- (b\i);}\end{scope}
 \begin{scope}[visible on=<6->]\foreach \i in {2,3,4}{\draw[flecha] (b\i) -- (c\i);}\end{scope}
\end{tikzpicture}
\elabpropia
\vspace{2pt}
\begin{clave}
Las promesas de esta etapa no quedan sueltas: los profetas las releen y Pablo las reúne en Romanos 11 (Estudio~8).
\end{clave}
\end{frame}

\begin{frame}{Isaías: del día de Madián a Isaí}
\begin{columns}[T,onlytextwidth]
\begin{column}{0.48\textwidth}
\centering{\heb{וְאָשִׁיבָה שֹׁפְטַיִךְ כְּבָרִאשֹׁנָה}}\par\smallskip
\begin{cita}Restauraré tus jueces como al principio [\ldots]; entonces te llamarán Ciudad de justicia, Ciudad fiel.\fuente{(Is 1:26)}\end{cita}
\footnotesize
\begin{itemize}
\item «Como en el día de Madián» (9:4 [9:3]); la peña de Oreb (10:26).
\item El niño «sobre el trono de David» (9:7 [9:6]).
\end{itemize}
\end{column}
\begin{column}{0.48\textwidth}
\centering{\heb{וְיָצָא חֹטֶר מִגֵּזַע יִשָׁי}}\par\smallskip
\begin{cita}Saldrá una vara del tronco de Isaí [\ldots] y reposará sobre él el Espíritu de YHWH.\fuente{(Is 11:1--2)}\end{cita}
\footnotesize
\begin{itemize}
\item «De Isaí», no «de David»: regreso a Belén (1~S 16).
\item Estandarte a los pueblos (11:10; Ro 15:12).
\end{itemize}
\end{column}
\end{columns}
\vspace{4pt}
\begin{notalam}[Filología: \heb{וְאָשִׁיבָה}]
Raíz \emph{shuv}, la del regreso: la restauración de Is 1:26 es una vuelta al ideal de los jueces que Dios levantaba. La misma raíz reaparece en Os 3:5.
\end{notalam}
\end{frame}

\begin{frame}{Jeremías, Oseas y Miqueas}
\begin{columns}[T,onlytextwidth]
\begin{column}{0.48\textwidth}
\begin{cita}[fontupper=\scriptsize\itshape]Levantaré a David un Renuevo justo [\ldots]. En sus días será salvo Judá, e Israel habitará confiado; y este será su nombre: YHWH, justicia nuestra.\fuente{(Jer 23:5--6)}\end{cita}
\scriptsize
\begin{itemize}
\item Silo como advertencia (Jer 7:12); «pastores según mi corazón» (3:15); «David su rey» (30:9).
\end{itemize}
\end{column}
\begin{column}{0.48\textwidth}
\begin{cita}[fontupper=\scriptsize\itshape]Muchos días estarán los hijos de Israel sin rey [\ldots], sin efod y sin terafines. Después volverán y buscarán a YHWH su Dios y a David su rey.\fuente{(Os 3:4--5)}\end{cita}
\scriptsize
\begin{itemize}
\item Efod y terafines: Jue 8:27; 17:5. «Días de Guibeá» (Os 10:9).
\item \textbf{Mi 5:2 [5:1]}: «tú, Belén Efrata [\ldots], de ti me saldrá el que será gobernante en Israel».
\end{itemize}
\end{column}
\end{columns}
\vspace{2pt}
\begin{notalam}[Precisión: Israel y Judá en los profetas]
En Jer 23:6, «Judá» es el reino del sur e «Israel» el reino del norte, ya desterrado: ambos bajo el rey davídico. En Oseas, «Israel» es el reino del norte, separado de la casa de David, que regresará a «David su rey».
\end{notalam}
\end{frame}

\begin{frame}{Pablo, benjaminita}
\fotolado{panel_ro11.jpg}{Elaboración propia: Ro 11:1 en escritura uncial (texto NA28).}
\begin{columns}[c,onlytextwidth]
\begin{column}{0.57\textwidth}
\begin{cita}¿Ha desechado Dios a su pueblo? En ninguna manera. Porque también yo soy israelita, de la descendencia de Abraham, de la tribu de Benjamín.\fuente{(Ro 11:1)}\end{cita}
\scriptsize
\begin{itemize}
\item \grc{ἀπώσατο}: el verbo de 1~S 12:22 (LXX).
\item \textbf{Benjamín}, casi extinguido (Jue 21), tribu de Saúl y de Saulo (Hch 13:9).
\item \textbf{Hch 13:22}: «he hallado a David [\ldots], varón conforme a mi corazón».
\item \textbf{Ro 1:3}: «del linaje de David según la carne».
\item \textbf{Ro 11:26} cita Is 59:20: el \heb{גּוֹאֵל} (\emph{go'el}), título de Booz.
\end{itemize}
\end{column}
\begin{column}{0.41\textwidth}\end{column}
\end{columns}
\end{frame}

\begin{frame}{Benjamín: de la casi extinción al apóstol}
\centering
\begin{tikzpicture}[x=1cm,y=1cm]
 \draw[line width=1.6pt,cafe,-{Latex[length=3mm]}] (0,0) -- (11.9,0);
 \foreach \x/\t/\b in {0.7/{Guibeá: Benjamín casi desaparece}/{Jue 19--21},
                       3.4/{Saúl, el benjaminita, libra a Jabes}/{1~S 11},
                       6.1/{Saúl rechazado; el reino «rasgado»}/{1~S 15:28},
                       8.8/{Saulo, «de la tribu de Benjamín»}/{Hch 13:9; Flp 3:5},
                       11.2/{«No ha desechado Dios a su pueblo»}/{Ro 11:1--2}}{
   \node[hito] at (\x,0) {};
   \node[above=5pt,font=\scriptsize\bfseries,text=ladrillo,text width=2.2cm,align=center] at (\x,0.05) {\t};
   \node[below=5pt,font=\scriptsize,text=sepia,text width=2.2cm,align=center] at (\x,-0.05) {\b};}
\end{tikzpicture}
\elabpropia
\vspace{8pt}
\begin{notalam}[Precisión]
Que sea un benjaminita quien pregunte si Dios desechó a su pueblo, y se presente como prueba de lo contrario, tiene un peso particular. Es una coincidencia que señala este estudio, no el apóstol.
\end{notalam}
\end{frame}

% ===================================================
\AtBeginSection[]{}
\section{Síntesis}
\begin{frame}{Síntesis}
\begin{tcolorbox}[colback=piedra,colframe=ladrillo,boxrule=1pt,arc=3pt,fontupper=\small]
\textbf{\textcolor{cafe}{De los Jueces a David:}} el pueblo llamado a ser reino de sacerdotes quiso ser «como todas las naciones» y se corrompió desde dentro, hasta convertirse en Sodoma en Guibeá. Dios, con todo, se compadeció, guardó un remanente (Débora, Rut, Ana, Samuel, Jonatán, David) y anunció un rey ungido salido de Belén. Sobre esa promesa descansa la esperanza de Isaías en el renuevo de Isaí, la de Jeremías en el Renuevo justo y la de Oseas en «David su rey», y de ella toma Pablo su respuesta: «No ha desechado Dios a su pueblo» (Ro 11:2).
\end{tcolorbox}
\scriptsize\renewcommand{\arraystretch}{1.25}
\begin{tabular}{>{\RaggedRight\arraybackslash}p{3.0cm}>{\RaggedRight\arraybackslash}p{8.6cm}}
\toprule
{\color{ladrillo}\bfseries Eje} & {\color{ladrillo}\bfseries Lo que deja esta etapa}\\\midrule
Vocación quebrada & Jue 17--21; 1~S 2:12--17; 8:5, 20\\
Remanente & Jue 5; Rt 1:16; 1~S 1--2; 7; 12:23; 17:46\\
Restauración & Jue 2:18; Rt 4:14--17; 1~S 2:10, 35; 12:22; 16:13\\
Proyección & Is 1:26; 11:1; Jer 23:5--6; Os 3:4--5; Mi 5:2 [5:1]; Ro 11:1--5\\
\bottomrule
\end{tabular}
\elabpropia
\end{frame}

\begin{frame}{Hacia el Estudio 4}
\begin{columns}[c,onlytextwidth]
\begin{column}{0.52\textwidth}
\begin{cita}Y vinieron los hombres de Judá y ungieron allí a David por rey sobre la casa de Judá.\fuente{(2~S 2:4)}\end{cita}
\small Un rey ungido, pero un reino incompleto: David reina en Hebrón sobre Judá, y la casa de Saúl conserva el resto de Israel.
\end{column}
\begin{column}{0.44\textwidth}
\centering
\begin{tikzpicture}[node distance=0.28cm]
 \node[bloque,text width=4.0cm] (h) {David en Hebrón, rey de Judá (2~S 2:4)};
 \begin{scope}[visible on=<2->]\node[bloqueoro,text width=4.0cm,below=of h] (j) {Jerusalén};\end{scope}
 \begin{scope}[visible on=<3->]\node[bloqueoro,text width=4.0cm,below=of j] (a) {El arca};\end{scope}
 \begin{scope}[visible on=<4->]\node[bloqueoro,text width=4.0cm,below=of a] (c) {La casa eterna (2~S 7)};\end{scope}
 \begin{scope}[visible on=<5->]\node[bloqueazul,text width=4.0cm,below=of c] (e) {Estudio 4: Salomón, el Primer Templo y la división entre Israel (norte) y Judá (sur)};\end{scope}
 \begin{scope}[visible on=<5->]\foreach \p/\q in {h/j,j/a,a/c,c/e}{\draw[flecha] (\p) -- (\q);}\end{scope}
\end{tikzpicture}
\elabpropia
\end{column}
\end{columns}
\end{frame}

\begin{frame}{Para la reflexión}
\small
\begin{enumerate}
\item El estribillo «cada uno hacía lo que bien le parecía» (Jue 21:25) retoma Dt 12:8. ¿Qué supone hoy vivir conforme a «lo recto ante los ojos de YHWH»?
\item Israel pidió un rey «como todas las naciones» (1~S 8:5). ¿Bajo qué formas aparece hoy el deseo de asemejarse al entorno?
\item Rut dijo «tu pueblo será mi pueblo, y tu Dios mi Dios». ¿Qué enseña su orden a los creyentes de las naciones (Ro 11:17--18)?
\item «Obedecer es mejor que el sacrificio» (1~S 15:22). ¿Cómo distinguir el culto como obediencia del culto como medio?
\item ¿Qué relación hay entre la preservación de Benjamín (Jue 21) y el remanente de Ro 11:1--5?
\end{enumerate}
\end{frame}

% ===================================================
\section{Glosario}
\begin{frame}{Glosario}
\begin{columns}[t,onlytextwidth]
\begin{column}{0.49\textwidth}\scriptsize
\textbf{Eben-ezer} (\heb{אֶבֶן הָעָזֶר}): «piedra de la ayuda» (1~S 7:12).\par\smallskip
\textbf{Go'el} (\heb{גֹּאֵל}): pariente redentor (Rt 4:14; Is 59:20).\par\smallskip
\textbf{Icabod} (\heb{אִי־כָבוֹד}): «¿dónde está la gloria?» (1~S 4:21).\par\smallskip
\textbf{Jésed} (\heb{חֶסֶד}): lealtad amorosa (Rt 1:8; 3:10).\par\smallskip
\textbf{Kejol ha-goyim} (\heb{כְּכָל־הַגּוֹיִם}): «como todas las naciones» (1~S 8:5).\par\smallskip
\textbf{Mashíaj} (\heb{מָשִׁיחַ}; \grc{χριστός}): ungido (1~S 2:10).
\end{column}
\begin{column}{0.49\textwidth}\scriptsize
\textbf{Naguid} (\heb{נָגִיד}): príncipe, designado (1~S 10:1).\par\smallskip
\textbf{Nazir} (\heb{נָזִיר}): consagrado por voto (Nm 6; Jue 13:5).\par\smallskip
\textbf{Seraním} (\heb{סְרָנִים}): príncipes filisteos (Jue 3:3).\par\smallskip
\textbf{Shofet} (\heb{שֹׁפֵט}): juez, libertador (Jue 2:16).\par\smallskip
\textbf{Tsémaj} (\heb{צֶמַח}): renuevo davídico (Jer 23:5).\par\smallskip
\textbf{Yashar be-einav} (\heb{הַיָּשָׁר בְּעֵינָיו}): «lo recto ante sus ojos» (Jue 21:25).
\end{column}
\end{columns}
\end{frame}

% ===================================================
\section{Referencias}
\begin{frame}[allowframebreaks]{Referencias}
\tiny
\refapa{\emph{Biblia Hebraica Stuttgartensia} (5.ª ed.). (1997). Deutsche Bibelgesellschaft.}
\refapa{Block, D.~I. (1999). \emph{Judges, Ruth} (New American Commentary, Vol. 6). Broadman \& Holman.}
\refapa{Bunimovitz, S., \& Lederman, Z. (2009). The archaeology of border communities: Renewed excavations at Tel Beth-Shemesh, Part 1: The Iron Age. \emph{Near Eastern Archaeology, 72}(3), 114--142.}
\refapa{Cohen, A. (Ed.). (1965). \emph{The minor tractates of the Talmud}. Soncino Press.}
\refapa{Cross, F.~M., \& Freedman, D.~N. (1997). \emph{Studies in ancient Yahwistic poetry} (2.ª ed.). Eerdmans.}
\refapa{Cross, F.~M., Parry, D.~W., Saley, R.~J., \& Ulrich, E. (2005). \emph{Qumran Cave 4.XII: 1--2 Samuel} (Discoveries in the Judaean Desert 17). Clarendon Press.}
\refapa{Danby, H. (Trad.). (1933). \emph{The Mishnah}. Oxford University Press.}
\refapa{Dothan, T. (1982). \emph{The Philistines and their material culture}. Yale University Press.}
\refapa{Epstein, I. (Ed.). (1935--1948). \emph{The Babylonian Talmud}. Soncino Press.}
\refapa{Faust, A. (2006). \emph{Israel's ethnogenesis: Settlement, interaction, expansion and resistance}. Equinox.}
\refapa{Feldman, M., Master, D.~M., Bianco, R.~A., Burri, M., Stockhammer, P.~W., Mittnik, A., Aja, A.~J., Jeong, C., \& Krause, J. (2019). Ancient DNA sheds light on the genetic origins of early Iron Age Philistines. \emph{Science Advances, 5}(7), eaax0061.}
\refapa{Finkelstein, I. (1988). \emph{The archaeology of the Israelite settlement}. Israel Exploration Society.}
\refapa{Finkelstein, I., Bunimovitz, S., \& Lederman, Z. (1993). \emph{Shiloh: The archaeology of a biblical site}. Institute of Archaeology, Tel Aviv University.}
\refapa{Finkelstein, I., \& Fantalkin, A. (2012). Khirbet Qeiyafa: An unsensational archaeological and historical interpretation. \emph{Tel Aviv, 39}(1), 38--63.}
\refapa{Freedman, H., \& Simon, M. (Eds.). (1939). \emph{Midrash Rabbah}. Soncino Press.}
\refapa{Garfinkel, Y., Kreimerman, I., \& Zilberg, P. (2016). \emph{Debating Khirbet Qeiyafa: A fortified city in Judah from the time of King David}. Israel Exploration Society.}
\refapa{Gitin, S., Dothan, T., \& Naveh, J. (1997). A royal dedicatory inscription from Ekron. \emph{Israel Exploration Journal, 47}(1--2), 1--16.}
\refapa{Guggenheimer, H.~W. (1998). \emph{Seder Olam: The rabbinic view of biblical chronology}. Jason Aronson.}
\refapa{Harrington, D.~J., \& Saldarini, A.~J. (1987). \emph{Targum Jonathan of the Former Prophets} (The Aramaic Bible, Vol. 10). Michael Glazier.}
\refapa{Hesse, B., \& Wapnish, P. (1997). Can pig remains be used for ethnic diagnosis in the ancient Near East? En N.~A. Silberman \& D.~B. Small (Eds.), \emph{The archaeology of Israel}. Sheffield Academic Press.}
\refapa{Josefo, F. (1987). \emph{The works of Josephus} (W. Whiston, Trad.). Hendrickson.}
\refapa{Levine, É. (1973). \emph{The Aramaic version of Ruth}. Biblical Institute Press.}
\refapa{Maeir, A.~M., Wimmer, S.~J., Zukerman, A., \& Demsky, A. (2008). A late Iron Age I/early Iron Age IIA Old Canaanite inscription from Tell eṣ-Ṣâfī/Gath, Israel. \emph{Bulletin of the American Schools of Oriental Research, 351}, 39--71.}
\refapa{Mazar, A. (1990). \emph{Archaeology of the land of the Bible, 10,000--586 B.C.E.} Doubleday.}
\refapa{McCarter, P.~K. (1980). \emph{I Samuel} (Anchor Bible, Vol. 8). Doubleday.}
\refapa{Moo, D.~J. (1996). \emph{The epistle to the Romans}. Eerdmans.}
\refapa{Neusner, J. (Trad.). (1982--1994). \emph{The Talmud of the Land of Israel}. University of Chicago Press.}
\refapa{\emph{Novum Testamentum Graece} (28.ª ed.). (2012). Deutsche Bibelgesellschaft.}
\refapa{ONE FOR ISRAEL Ministry. (2025, 3 de julio). \emph{What does «all Israel will be saved» mean in Romans 11?} \url{https://www.oneforisrael.org/what-does-all-israel-will-be-saved-mean-in-romans-11/}}
\refapa{Rahlfs, A., \& Hanhart, R. (Eds.). (2006). \emph{Septuaginta} (ed. rev.). Deutsche Bibelgesellschaft.}
\refapa{Sacks, J. (Trad.). (2009). \emph{The Koren Siddur}. Koren Publishers.}
\refapa{Stager, L.~E. (1985). The archaeology of the family in ancient Israel. \emph{Bulletin of the American Schools of Oriental Research, 260}, 1--35.}
\end{frame}
\end{document}
